\subsection{Variational principles for the spectrum of positive operators}

In this number, E is a nonzero complex Hilbert space.

PROPOSITION 15. --- Let u be a self-adjoint partial operator on E. Suppose that u is bounded below by \(c \in \mathbf{R}\) (def. 7 of IV, p. 237). We have

\[
\inf (\mathrm{Sp} (u)) = \inf _ {x \in \mathrm{dom} (u) - \{0 \}} \frac {\langle x \mid u (x) \rangle}{\| x \| ^ {2}} \in [ c, + \infty [.
\]

Let m be the right-hand side of the equality to be proved. Let \(x \in E\) and \(\mu_{x}\) be the spectral measure of x relative to u. We have

\[
\begin{array}{r l} \langle x \mid u (x) \rangle & = \int_ {\mathrm{Sp} (u)} t d \mu_ {x} (t) \\ & \geqslant \inf (\mathrm{Sp} (u)) \int_ {\mathrm{Sp} (u)} d \mu_ {x} (t) = \inf (\mathrm{Sp} (u)) \| x \| ^ {2}, \end{array}
\]

(formula (6), p. 271), hence \(\inf(\mathrm{Sp}(u)) \leqslant m\) . Conversely, the partial operator \(u - m\) is positive, therefore \(\inf(\mathrm{Sp}(u)) - m = \inf(\mathrm{Sp}(u - m \cdot 1_{\mathrm{E}})) \geqslant 0\) (prop. 17 of IV, p. 248).

Proposition 9 of I, p. 139 corresponds to the particular case of this proposition when u is a Hermitian element of \(\mathcal{L}(\mathrm{E})\) , which is then necessarily bounded below. In this case, we also have

\[
\sup (\operatorname{Sp} (u)) = \sup _ {x \in \operatorname{dom} (u) - \{0 \}} \frac {\langle x \mid u (x) \rangle}{\| x \| ^ {2}}
\]

(loc. cit.); if u does not belong to \(\mathcal{L}(\mathrm{E})\) then this supremum is \(+\infty\) .

DEFINITION 10. --- Let u be a normal partial operator on E. A complex number \(\lambda \in \mathrm{Sp}(u)\) belongs to the sensitive spectrum of \(u\) if \(\lambda\) is isolated in \(\mathrm{Sp}(u)\) and if \(\lambda\) is an eigenvalue of finite multiplicity.

We denote by \(\mathrm{Sp}_{\mathrm{s}}(u)\) the sensitive spectrum of u. Its complement in \(\mathrm{Sp}(u)\) is called the essential spectrum of u and denoted \(\mathrm{Sp}_{\mathrm{e}}(u)\) .

The essential spectrum of a normal partial operator u on E is closed in C, since \(\mathrm{Sp}(u)\) is closed in C and the complement of the essential spectrum is open in \(\mathrm{Sp}(u)\). The sensitive spectrum of u is not necessarily closed in C (Exercise 36 of IV, p. 369).

Lemma 9. --- Let E be a complex Hilbert space and let u be a normal partial operator on E. Let \(\lambda \in \mathrm{Sp}(u)\) . We have \(\lambda \in \mathrm{Sp}_{\mathrm{s}}(u)\) if and only if there exists an open neighborhood V of \(\lambda\) in C such that the spectral projector \(p_{\mathrm{V}}\) of u defined by V has finite rank.

Let \(\lambda \in \mathrm{Sp}_{\mathrm{s}}(u)\). There exists an open neighborhood \(V\) of \(\lambda\) in \(\mathbf{C}\) such that \(\mathrm{Sp}(u) \cap V = \{\lambda\}\) ; the spectral projector \(p_V = p_{\{\lambda\}}\) is then of finite rank (cor. of prop. 9 of IV, p. 278).

Conversely, suppose that there exists an open neighborhood V of \(\lambda\) such that the spectral projector \(p_{V}\) has finite rank \(n \in N\). By the corollary of Proposition 10 of IV, p. 279, the intersection \(V \cap \mathrm{Sp}(u)\) contains at most n elements, therefore \(\lambda\) is isolated in \(\mathrm{Sp}(u)\). This is therefore an eigenvalue of u of spectral multiplicity at most n, by the corollary of prop. 9 of IV, p. 278.

In the rest of this number, we assume that E is infinite-dimensional; the analogues of the results below when E is finite-dimensional are deduced from No. 4 of IV, p. 153.

Let u be a self-adjoint operator bounded below on E. Let c be a lower bound of u; the spectrum of u is contained in \([c, +\infty[\) (prop. 15). Suppose that the essential spectrum of u is not empty. We then denote by \(\varrho_{e}\) the lower bound of the essential spectrum of u, so that \(\varrho_{e} \geqslant c\) and \(\varrho_{e}\) is an element of the spectrum of \(u\) . We denote \(\mathrm{Sp}_{\mathrm{h}}(u) = \mathrm{Sp}(u) \cap [\varrho_{\mathrm{e}}, +\infty[\); it is a closed subset of \(\mathrm{Sp}(u)\) , hence also of \(\mathbf{C}\) , such that \(\inf (\mathrm{Sp}_{\mathrm{h}}(u)) = \varrho_{\mathrm{e}}\) ; it is called the upper part of the spectrum of \(u\) . If the essential spectrum of \(u\) is empty, we denote \(\mathrm{Sp}_{\mathrm{h}}(u) = \varnothing\) and \(\varrho_{\mathrm{e}} = +\infty\) .

The intersection \(\mathrm{Sp}_{\mathrm{b}}(u) = \mathrm{Sp}(u) \cap [0, \varrho_{\mathrm{e}}[\) is contained in the sensitive spectrum of \(u\) and its elements are therefore isolated eigenvalues of finite multiplicity; this is called the lower part of the spectrum of \(u\) . Let \(\mathrm{E}_{\mathrm{b}}\) be the closed subspace of E generated by the eigenspaces corresponding to the eigenvalues \(\lambda \in \mathrm{Sp}_{\mathrm{b}}(u)\) . By the cor. of prop. 9 of IV, p. 278 and prop. 8 of IV, p. 278, the space \(\mathrm{E}_{\mathrm{b}}\) is the image of the spectral projector \(p_{\mathrm{Sp}_{\mathrm{b}}(u)}\) defined by \(\mathrm{Sp}_{\mathrm{b}}(u)\) . Since \(\mathrm{Sp}(u)\) is the disjoint union of \(\mathrm{Sp}_{\mathrm{b}}(u)\) and \(\mathrm{Sp}_{\mathrm{h}}(u)\) , the orthogonal complement \(\mathrm{E}_{\mathrm{h}}\) of \(\mathrm{E}_{\mathrm{b}}\) is the image of the spectral projector \(p_{\mathrm{Sp}_{\mathrm{h}}(u)}\) defined by \(\mathrm{Sp}_{\mathrm{h}}(u)\) . If \(\mathrm{Sp}_{\mathrm{b}}(u)\) is finite, then the space \(\mathrm{E}_{\mathrm{b}}\) is finite-dimensional; the space \(\mathrm{E}_{\mathrm{h}}\) is then nonzero, since E is assumed to be infinite-dimensional.

Lemma 10. --- One has

\[
\langle x \mid u (x) \rangle \geqslant \varrho_ {\mathrm{e}} \| x \| ^ {2}.\tag{18}
\]

for all \(x \in \mathrm{E}_{\mathrm{h}} \cap \operatorname{dom}(u)\) .

If \(x \in \mathrm{E}_{\mathrm{h}} \cap \mathrm{dom}(u)\) , then the support of the spectral measure \(\mu_{x}\) of \(x\) relative to \(u\) is contained in the interval \([\varrho_{\mathrm{e}}, +\infty[\) (prop. 9, \(a\) ) of IV, p. 278), hence

\[
\langle x \mid u (x) \rangle = \int_ {\mathbf {R}} t d \mu_ {x} (t) \geqslant \varrho_ {\mathrm{e}} \| x \| ^ {2}.
\]

Lemma 11. --- Suppose that \(E_{b}\) is finite-dimensional. Then, for every real number \(\varepsilon > 0\), the image of the spectral projector of u defined by \([\varrho_{e}, \varrho_{e} + \varepsilon]\) is infinite-dimensional.

If \(E_{b}\) is finite-dimensional, then the essential spectrum of u is nonempty, hence \(\varrho_{e}\) is finite and belongs to \(\mathrm{Sp}_{\mathrm{e}}(u)\) . Moreover, the lower spectrum \(\mathrm{Sp}_{\mathrm{b}}(u)\) is finite, so there exists \(\delta > 0\) such that \([\varrho_{e} - \delta, \varrho_{e} + \delta] \cap \mathrm{Sp}(u) \subset [\varrho_{e}, \varrho_{e} + \delta]\) . The assertion then follows from Lemma 9.

PROPOSITION 16. --- The set \(\mathrm{Sp}_{\mathrm{b}}(u) \subset [c, \varrho_{\mathrm{e}}[\) is countable, discrete, and closed in \([0, \varrho_{\mathrm{e}}[\) . It is the set of values of a strictly increasing sequence \((\nu_{k})_{0 \leqslant k < \mathrm{Card}(\mathrm{Sp}_{\mathrm{b}}(u))}\) of positive real numbers; if \(\mathrm{Sp}_{\mathrm{b}}(u)\) is infinite, then the sequence \((\nu_{k})\) converges to \(\varrho_{\mathrm{e}}\) .

Let \(\mathrm{T} = \mathrm{Sp}_{\mathrm{b}}(u) \cap [0, \varrho_{\mathrm{e}}[\) . It is a closed and discrete subset of \([c, \varrho_{\mathrm{e}}[\) by definition. For every integer \(i \geqslant 1\) , the set \(\mathrm{Sp}_{\mathrm{b}}(u) \cap [c, \varrho_{\mathrm{e}} - 1/i]\) is compact and discrete, hence finite. Since T is the union of these sets for \(i \geqslant 1\) , the set T is countable.

This concludes the proof if \(\mathrm{Sp}_{\mathrm{b}}(u)\) is finite. If \(\mathrm{Sp}_{\mathrm{b}}(u)\) is infinite, one concludes by applying Lemma 2 of III, p. 90 to the image S of T under the map \(\lambda \mapsto \varrho_{e} - \lambda\) .

For every integer \(k\) such that \(0 \leqslant k < \text{Card}(\text{Sp}_b(u))\) , we denote by \(n_k \geqslant 1\) the multiplicity of the eigenvalue \(\nu_k\) of \(u\) . We denote \(\overline{\mathbf{N}} = \mathbf{N} \cup \{+\infty\} \subset \overline{\mathbf{R}}\) . Let \(\mathbf{M} \in \overline{\mathbf{N}}\) be the sum of the multiplicities \(n_k\) . This is the Hilbertian dimension of \(\mathbf{E}_b\) , if one agrees to say that a space of infinite Hilbertian dimension has dimension \(+\infty \in \overline{\mathbf{N}}\) .

For every integer n such that \(0 \leqslant n < M\) , one defines \(\lambda_{n}(u) = \nu_{k}\) , where \(k \geqslant 0\) is the unique integer such that

\[
n _ {0} + \dots + n _ {k - 1} \leqslant n <   n _ {0} + \dots + n _ {k}.
\]

If \(n \in \mathbf{N}\) satisfies \(n \geqslant \mathrm{M}\), we set \(\lambda_n(u) = \varrho_{\mathrm{e}}\) . This case can occur only if \(\mathrm{Sp}_{\mathrm{b}}(u)\) is finite, in which case \(\varrho_{\mathrm{e}}\) is finite, since E is assumed to be of infinite dimension.

By construction, the sequence \((\lambda_{n}(u))_{n\in\mathbf{N}}\) is increasing and, for every element \(\lambda\) of \(\mathrm{Sp}_{\mathrm{b}}(u)\) , the number of integers n such that \(\lambda_{n}(u)=\lambda\) is equal to the multiplicity of \(\lambda\) as an eigenvalue of u. The sequence \((\lambda_{n}(u))_{n\in\mathbf{N}}\) tends to \(+\infty\) if and only if the essential spectrum of u is empty.

For every \(n \in N\) , we denote by \(F_{n}\) (resp. \(F_{n}^{u}\) ) the set of vector subspaces \(F \subset E\) of dimension n (resp. the set of vector subspaces \(F \subset \text{dom}(u)\) of dimension n).

Let \(n \in \mathbf{N}\) such that \(n < \mathrm{M}\) and \(\mathrm{F} \in \mathcal{F}_n^u\) . One says that \(\mathrm{F}\) is adapted to \(u\) if \(\mathrm{F}\) admits an orthonormal basis \((f_i)_{0 \leqslant i \leqslant n-1}\) such that \(u(f_i) = \lambda_i(u)f_i\) for \(0 \leqslant i \leqslant n-1\) .

Let \(F \in F_{n}^{u}\) adapted to u. The space F is contained in \(E_{b}\) ; it is generated by eigenvectors of u for eigenvalues \(\lambda \leqslant \lambda_{n-1}(u)\) , and it contains the eigenspaces corresponding to the eigenvalues \(\lambda < \lambda_{n-1}(u)\) . Moreover, for every universally measurable subset A of \(\mathrm{Sp}(u)\) , the space F is stable under the spectral projector \(p_{A}\) defined by A (prop. 9, c) of IV, p. 278).

Lemma 12. --- Let \(F \in F_{n}^{u}\) be a subspace adapted to u. Every eigenvector of u belonging to the space \(F^{\circ} \cap E_{b}\) has an eigenvalue \(\geqslant \lambda_{n}(u)\) .

Let \(\ell\) be such that \(\lambda_{n-1}(u) = \nu_{\ell}\) . Let \(x \in \mathrm{F}^{\circ} \cap \mathrm{E}_{\mathrm{b}}\) be an eigenvector of \(u\) for an eigenvalue \(\lambda\) and let \(k < \operatorname{Card}(\operatorname{Sp}_{\mathrm{b}}(u))\) be such that \(\lambda = \nu_{k}\) .

The condition \(k < \ell\) is impossible, because \(\nu_k < \nu_\ell\) , and the space F would then contain the eigenspace for the eigenvalue \(\nu_k\) , contradicting the fact that \(x\) is orthogonal to F. If \(k = \ell\) , then \(x\) is an eigenvector for the eigenvalue \(\lambda_{n-1}(u)\) ; since \(x \in \mathrm{F}^\circ\) , the multiplicity \(n_k\) is strictly greater than the number of integers \(i < n\) such that \(\lambda_i(u) = \nu_k\) , which implies that \(\lambda_n(u) = \lambda_{n-1}(u)\) . Finally, if \(k > \ell\) , we have \(\lambda = \nu_k > \nu_\ell = \lambda_{n-1}(u)\) , hence \(\lambda \geqslant \lambda_n(u)\) .

For every subspace F of E, we denote by

\[
\widetilde{r}_{\mathrm{F}}(u) = \inf_{\substack{x\in \operatorname{dom}(u)\cap \mathrm{F}^{\circ}\\ x\neq 0}}\frac{\langle x\mid u(x)\rangle}{\|x\|^{2}},
\]

\[
\widetilde{\mathrm{R}}_{\mathrm{F}}(u) = \sup_{\substack{x\in \operatorname{dom}(u)\cap \mathrm{F}\\ x\neq 0}}\frac{\langle x\mid u(x)\rangle}{\|x\|^ {2}}.
\]

PROPOSITION 17. --- a) For every integer \(n \in \mathbf{N}\) , one has

\[
\lambda_ {n} (u) = \sup _ {\mathrm{F} \in \mathcal {F} _ {n}} \widetilde {r} _ {\mathrm{F}} (u) = \inf _ {\mathrm{F} \in \mathcal {F} _ {n + 1} ^ {u}} \widetilde {\mathrm{R}} _ {\mathrm{F}} (u);
\]

\begin{enumerate}
\def\labelenumi{\alph{enumi})}
\setcounter{enumi}{1}
\item
  For every integer \(n < \mathrm{M}\) and for every subspace \(\mathrm{F} \in \mathcal{F}_n^u\) adapted to \(u\) , we have \(\lambda_n(u) = \widetilde{r}_{\mathrm{F}}(u)\) ;
\item
  For every integer \(n < \mathrm{M}\) and for every subspace \(\mathrm{F} \in \mathcal{F}_{n+1}^{u}\) adapted to \(u\) , we have \(\lambda_{n}(u) = \widetilde{\mathrm{R}}_{\mathrm{F}}(u)\) .
\end{enumerate}

There exists an orthonormal basis \((e_{n})_{0\leqslant n<M}\) of the space \(E_{b}\) such that \(e_{n}\in\operatorname{dom}(u)\) and \(u(e_{n})=\lambda_{n}(u)e_{n}\) for every n such that \(0\leqslant n<M\) . For every \(x\in E_{b}\cap\operatorname{dom}(u)\) , we therefore have

\[
\langle x \mid u (x) \rangle = \sum_ {0 \leqslant n <   M} \lambda_ {n} (u) | \langle e _ {n} \mid x \rangle | ^ {2}.
\]

For every integer \(n\) such that \(1 \leqslant n < M + 1\) , let \(F_n\) be the subspace of dimension \(n\) of \(E_b\) generated by \((e_0, \ldots, e_{n-1})\) . We have \(F_n \subset \text{dom}(u)\) and \(F_n\) is adapted to \(u\) .

Let \(n \in \mathbf{N}\) and \(\mathrm{F} \in \mathcal{F}_n\) . Let us prove that \(\widetilde{r}_{\mathrm{F}}(u) \leqslant \lambda_n(u)\) and, consequently, that

\[
\sup _ {\mathrm{F} \in \mathcal {F} _ {n}} \widetilde {r} _ {\mathrm{F}} (u) \leqslant \lambda_ {n} (u).\tag{19}
\]

If \(0 \leqslant n < M\) , in particular if M is infinite, then the restriction to \(\mathrm{F}_{n+1}\) of the orthogonal projection of E onto F is not injective, hence there exists \(x \neq 0\) in \(F_{n+1}\) orthogonal to F. Since

\[
\langle x \mid u (x) \rangle = \sum_ {0 \leqslant i \leqslant n} \lambda_ {i} (u) | \langle e _ {i} \mid x \rangle | ^ {2} = \lambda_ {n} (u) | \langle e _ {n} \mid x \rangle | ^ {2} \leqslant \lambda_ {n} (u) \| x \| ^ {2},
\]

we have \(\widetilde{r}_{\mathrm{F}}(u) \leqslant \lambda_n(u)\) .

If \(M \in N\) and \(n \geqslant M\) , then for every real number \(\varepsilon > 0\) , there exists x of norm 1 in \(\text{dom}(u)\) which is orthogonal to F and whose spectral measure \(\mu_{x}\) has support contained in \([\varrho_{e}, \varrho_{e} + \varepsilon]\) (lemma 11 and prop. 9, a) of IV, p. 278), whence

\[
\widetilde {r} _ {\mathrm{F}} (u) \leqslant \langle x \mid u (x) \rangle = \int_ {\operatorname{Sp} (u)} t d \mu_ {x} (t) \leqslant \varrho_ {\mathrm{e}} + \varepsilon = \lambda_ {n} (u) + \varepsilon
\]

by definition. Since \(\varepsilon > 0\) is arbitrary, we therefore have \(\widetilde{r}_{\mathrm{F}}(u) \leqslant \lambda_n(u)\) . Inequality (19) is therefore established.

Let \(n \in \mathbf{N}\) and \(\mathrm{F} \in \mathcal{F}_{n+1}^{u}\) . Let us prove that \(\widetilde{\mathrm{R}}_{\mathrm{F}}(u) \geqslant \lambda_{n}(u)\) and, consequently, that

\[
\inf _ {\mathrm{F} \in \mathcal {F} _ {n + 1} ^ {u}} \widetilde {\mathrm{R}} _ {\mathrm{F}} (u) \geqslant \lambda_ {n} (u).\tag{20}
\]

If \(0 \leqslant n < M\) , observe that the restriction to F of the orthogonal projector onto \(\mathrm{F}_n\) is not injective, hence there exists a vector \(x \neq 0\) in F orthogonal to \(\mathrm{F}_n\) . Set \(x_{\mathrm{b}} = p_{\mathrm{Sp}_{\mathrm{b}}(u)}(x)\) and \(x_{\mathrm{h}} = p_{\mathrm{Sp}_{\mathrm{h}}(u)}(x)\) . We therefore have \(x = x_{\mathrm{b}} + x_{\mathrm{h}}\) . The elements \(x_{\mathrm{b}}\) and \(x_{\mathrm{h}}\) belong to the domain of \(u\) (prop. 9, c) of IV, p. 278) and are orthogonal to \(\mathrm{F}_n\) . We have the lower bound

\[
\langle x _ {\mathrm{b}} \mid u (x _ {\mathrm{b}}) \rangle = \sum_ {i \geqslant n} \lambda_ {i} (u) | \langle e _ {i} \mid x _ {\mathrm{b}} \rangle | ^ {2} \geqslant \lambda_ {n} (u) \| x _ {\mathrm{b}} \| ^ {2},
\]

and \(\langle x_{\mathrm{h}}|u(x_{\mathrm{h}})\rangle \geqslant \varrho_{\mathrm{e}}\| x_{\mathrm{h}}\|^2\) (formula (18)), whence

\[
\begin{array}{r l} & {\langle x \mid u (x) \rangle = \langle x _ {\mathrm{b}} \mid u (x _ {\mathrm{b}}) \rangle + \langle x _ {\mathrm{h}} \mid u (x _ {\mathrm{h}}) \rangle} \\ & {\qquad \geqslant \lambda_ {n} (u) \| x _ {\mathrm{b}} \| ^ {2} + \varrho_ {\mathrm{e}} \| x _ {\mathrm{h}} \| ^ {2} \geqslant \lambda_ {n} (u) \| x \| ^ {2}.} \end{array}
\]

If M is finite and \(n \geqslant M\) , there exists \(x \neq 0\) in F orthogonal to \(E_{b}\) , hence \(x \in E_{h}\) , and

\[
\langle x \mid u (x) \rangle \geqslant \varrho_ {\mathrm{e}} \| x \| ^ {2} = \lambda_ {n} (u) \| x \| ^ {2}
\]

by (18). Inequality (20) is therefore proved.

We shall now prove the assertions \(b\) ) and \(c\) ), which imply that inequalities (19) and (20) are equalities when \(n < M\) .

Let us prove assertion \(b\) ). Let \(\mathrm{F} \in \mathcal{F}_n^u\) be a space adapted to \(u\) . We have the Hilbert sum

\[
\mathrm{F}^{\circ} = \mathrm{E}_{\mathrm{h}}\oplus \bigoplus_{\substack{\lambda \in \operatorname{Sp}_{\mathrm{b}}(u)\\ \lambda >\lambda_{n - 1}(u)}}\operatorname{Ker}(u - \lambda \cdot 1_{\mathrm{E}})\oplus (\mathrm{F}^{\circ}\cap \operatorname{Ker}(u - \lambda_{n - 1}(u)\cdot 1_{\mathrm{E}})).
\]

Let \(x \in \mathrm{F}^{\circ} - \{0\}\) . Let us write \(x = x_{\mathrm{h}} + y + z\) , where \(x_{\mathrm{h}} \in \mathrm{E}_{\mathrm{h}}\) and \(y\) (resp. \(z\) ) belongs to the second (resp. third) space of the above decomposition. Using (18) again and the fact that every eigenvalue \(\lambda > \lambda_{n-1}(u)\) of \(u\) is \(\geqslant \lambda_n(u)\) , we obtain

\[
\begin{array}{r l} & {\langle x \mid u (x) \rangle = \langle x _ {\mathrm{h}} \mid u (x _ {\mathrm{h}}) \rangle + \langle y \mid u (y) \rangle + \langle z \mid u (z) \rangle} \\ & {\qquad \geqslant \varrho_ {\mathrm{e}} \| x _ {\mathrm{h}} \| ^ {2} + \lambda_ {n} (u) \| y \| ^ {2} + \lambda_ {n - 1} (u) \| z \| ^ {2}.} \end{array}
\]

If \(z \neq 0\) , then by Lemma 12, we have \(\lambda_n(u) = \lambda_{n-1}(u)\) . We therefore deduce that \(\langle x | u(x) \rangle \geqslant \lambda_n(u) \| x \|^2\) , whence \(\widetilde{r}_{\mathrm{F}}(u) \geqslant \lambda_n(u)\) . Combined with (19), this implies the assertion \(b\) ).

Let us prove assertion \(c\) ). Let \(\mathrm{F} \in \mathcal{F}_{n+1}^{u}\) be a subspace adapted to \(u\) . Let \((f_{i})_{0 \leqslant i \leqslant n}\) be an orthonormal basis of \(\mathrm{F}\) such that \(u(f_{i}) = \lambda_{i}(u)f_{i}\) for \(0 \leqslant i \leqslant n\) . For every \(x \in \mathrm{F}\) , we have

\[
\langle x \mid u (x) \rangle = \sum_ {0 \leqslant i \leqslant n} \lambda_ {i} (u) | \langle f _ {i} \mid x \rangle | ^ {2} \leqslant \lambda_ {n} (u) \| x \| ^ {2},
\]

with equality if \(x = f_{n}\) , hence \(\widetilde{\mathrm{R}}_{\mathrm{F}}(u) = \lambda_{n}(u)\) .

Let us finally prove that (19) and (20) are equalities when \(n \geqslant M\) . In this case, M is finite, hence \(E_{b}\) is contained in the domain of u; moreover, we have \(\lambda_{n}(u) = \varrho_{\mathrm{e}}\) by definition.

There exists \(F \in F_{n}\) containing \(E_{b}\) . Every element \(x \neq 0\) of \(\text{dom}(u)\) orthogonal to F is therefore orthogonal to \(E_{b}\) , and satisfies

\[
\frac {\langle x \mid u (x) \rangle}{\| x \| ^ {2}} \geqslant \varrho_ {\mathrm{e}}
\]

(formula (18)), hence \(\widetilde{r}_{\mathrm{F}}(u) \geqslant \varrho_{\mathrm{e}} = \lambda_n(u)\) , and consequently

\[
\sup _ {\mathrm{F} \in \mathcal {F} _ {n}} \widetilde {r} _ {\mathrm{F}} (u) \geqslant \lambda_ {n} (u).
\]

Let \(\varepsilon > 0\) . Since \(\mathrm{E_b}\) is finite-dimensional, there exists by Lemma 11 an orthonormal family \((x_i)_{i \in \mathrm{I}}\) of elements of E such that the subspace F generated by \(\mathrm{E_b}\) and \((x_i)_{i \in \mathrm{I}}\) has dimension \(n + 1\) and is contained in \(\operatorname{dom}(u)\) , and such that \(\langle x_{i} \mid u(x_{i})\rangle \leqslant \varrho_{\mathrm{e}} + \varepsilon\) for all \(i \in I\) . We then have \(\widetilde{\mathrm{R}}_{\mathrm{F}}(u) \leqslant \varrho_{\mathrm{e}} + \varepsilon\) . Since \(\varepsilon > 0\) is arbitrary, we conclude that

\[
\inf _ {\mathrm{F} \in \mathcal {F} _ {n + 1} ^ {u}} \widetilde {\mathrm{R}} _ {\mathrm{F}} (u) \leqslant \varrho_ {\mathrm{e}} = \lambda_ {n} (u).
\]

The proposition is proved.

